% Authoritative exercise source. Preserve all derivations and diagrams.

\ExerciseStart{G2.10}{Ricavare autovalori, autovettori e proiettori}{Consolidamento}

\BookHeading{Problema, dati e obiettivo}

Per \(A=\begin{pmatrix}2&1\\1&2\end{pmatrix}\), trovare autovalori, autovettori normalizzati e proiettori. Ricostruire A, calcolare \(A^2\) e applicare \(e^{-itA}\) a \(e_1\) per \(t=\pi/2\). Il parametro t è adimensionale.

\ExTheory{\ExRef{c2-spettro}{2.12}; \ExRef{c2-base-propria}{2.13}; \ExRef{c2-proiettori}{2.14}; \ExRef{c2-funzioni-matrice}{2.15}}{Un autovettore soddisfa \(Au=\lambda u\), con \(u\ne0\). Gli autovalori risolvono \(\det(A-\lambda I)=0\). Per una base propria ortonormale, \(P_j=u_ju_j^\dagger\), \(A=\sum_j\lambda_jP_j\) e \(f(A)=\sum_jf(\lambda_j)P_j\).}

\ExPrediction{La matrice tratta simmetricamente le due coordinate. Le diagonali del piano sono quindi candidate naturali ad autodirezioni. La componente su ciascuna diagonale subirà un fattore indipendente.}

\ExFigure{\begin{tikzpicture}[>=Stealth,exlabel]
\draw[exaxis] (-0.35,0)--(1.55,0) node[right] {$x_1$};\draw[exaxis] (0,-1.35)--(0,1.55) node[above] {$x_2$};\draw[->,Accent,line width=1.1pt] (0,0)--(1,1) node[anchor=south west] {$ u_+\sqrt2 $};\draw[->,Secondary,line width=1.1pt] (0,0)--(1,-1) node[anchor=north west] {$ u_-\sqrt2 $};
\end{tikzpicture}}{Le due diagonali sono autodirezioni ortogonali.}


\begin{ExTasks}

\item Sviluppare il polinomio caratteristico.

\item Risolvere i due sistemi degli autovettori.

\item Normalizzare e verificare l’ortogonalità.

\item Costruire i due prodotti esterni.

\item Ricostruire A e \(A^2\).

\item Calcolare le due fasi e la loro azione su \(e_1\).

\end{ExTasks}

\ExSolution

\ExStep{Trovare gli autovalori}

\ExEq{\det(A-\lambda I)&=(2-\lambda)^2-1\\&=\lambda^2-4\lambda+3=(\lambda-3)(\lambda-1).}
Gli autovalori sono \(\lambda_+=3\) e \(\lambda_-=1\).

\ExStep{Risolvere i sistemi}

Per \(\lambda=3\), \(-u_1+u_2=0\), quindi \(u_2=u_1\). Per \(\lambda=1\), \(u_1+u_2=0\), quindi \(u_2=-u_1\). Si possono scegliere \((1,1)^T\) e \((1,-1)^T\).

\ExStep{Normalizzare e controllare}

Entrambe le norme sono \(\sqrt{1+1}=\sqrt2\):
\ExEq{u_+&=\frac1{\sqrt2}\colv{1\\1},\quad u_-=\frac1{\sqrt2}\colv{1\\-1},\\u_+^\dagger u_-&=(1-1)/2=0.}
Direttamente, \(Au_+=(3,3)^T/\sqrt2=3u_+\) e \(Au_-=(1,-1)^T/\sqrt2=u_-\).

\ExStep{Costruire i proiettori}

\ExEq{P_+&=\tfrac12\colv{1\\1}(1,1)=\tfrac12\begin{pmatrix}1&1\\1&1\end{pmatrix},\\P_-&=\tfrac12\colv{1\\-1}(1,-1)=\tfrac12\begin{pmatrix}1&-1\\-1&1\end{pmatrix}.}
La somma è I; il prodotto è zero perché \(u_+^\dagger u_-=0\).

\ExStep{Ricostruire l’operatore}

\ExEq{3P_++P_-&=\tfrac12\begin{pmatrix}3+1&3-1\\3-1&3+1\end{pmatrix}
=\begin{pmatrix}2&1\\1&2\end{pmatrix}=A.}

\ExStep{Calcolare il quadrato}

\ExEq{A^2&=9P_++P_-=\tfrac12\begin{pmatrix}10&8\\8&10\end{pmatrix}
=\begin{pmatrix}5&4\\4&5\end{pmatrix}.}
La moltiplicazione diretta produce elementi \(2\cdot2+1\cdot1=5\) e \(2\cdot1+1\cdot2=4\), con gli altri due ottenuti analogamente dalle rispettive righe.

\ExStep{Assegnare le fasi alle autodirezioni}

\ExEq{e^{-itA}&=e^{-3it}P_++e^{-it}P_-,\\e^{-3i\pi/2}&=i,\qquad e^{-i\pi/2}=-i,\\e^{-i\pi A/2}&=i(P_+-P_-)=i\begin{pmatrix}0&1\\1&0\end{pmatrix}.}

\ExStep{Applicare e interpretare}

\ExEq{e^{-i\pi A/2}e_1&=i\colv{0\\1}.}
La norma finale è uno. Se A rappresentasse un Hamiltoniano diviso per una scala energetica e t fosse il tempo corrispondentemente riscalato, l’esponenziale descriverebbe l’evoluzione. Le due componenti proprie accumulano fasi differenti prima di ricombinarsi.

\par\noindent\begin{minipage}{\linewidth}
\ExFigure{\begin{tikzpicture}[>=Stealth,exlabel]
\node[exbox] (n0) at (0.0,0.0) {$e_1=(u_++u_-)/\sqrt2$};\node[exbox] (n1) at (3.6,0.0) {fattori $i,-i$};\draw[exother] (n0)--(n1);\node[exbox] (n2) at (7.2,0.0) {$i e_2$};\draw[exother] (n1)--(n2);
\end{tikzpicture}}{Scomposizione, evoluzione delle fasi e ricomposizione del vettore.}


\ExInterpret{La base propria separa l’azione dell’operatore in trasformazioni scalari. I proiettori riportano ciascuna componente nello spazio originario, permettendo di verificare il risultato senza cambiare il significato geometrico del vettore.}

\ExLab{Nelle sezioni 2.12–2.15 confrontare autovalori, base propria e proiettori della matrice assegnata. Nello strumento matriciale verificare \(A^2\) per moltiplicazione diretta. Le figure delle autodirezioni devono conservare la stessa convenzione di base.}


\end{minipage}
